% -----------------------------------------------------------
\section{1854}
% -----------------------------------------------------------

	% ----------------------
	\subsection{Brigham Young}
	% ----------------------
		\label{EMS-1854-BY}

		% ----------------------
		\subsubsection{Introduction to 1854}
		% ----------------------
		
			sdf
			
		% -----
		\subsubsection{Instruction, 12 February 1854}
		% -----
			\label{EMS-1854-BY-12-FEB-1854}
			% \label{EMS-18XX-BY-DAY-3LETTERMONTH-YEAR}
			
			\begin{description}
					\item[Print Sources]
				\end{description}

					\begin{tabular}{p{0.5in}p{1in}p{0.5in}p{1in}}
					CDBY	& 		&  		&  \\
					\end{tabular}
				
				\begin{description}
					\item[Online Access] \href{https://catalog.churchofjesuschrist.org/assets/1e4f947a-0644-47b0-b058-5dd134a82a00/0/0}{Here}
					% \item[Analysis] \hyperref[XX]{Here}
				\end{description}
				
				\begin{description}
					\item[Background]
				\end{description}
				
					\phantom{.}
				
				\begin{description}
					\item[Original Source]
				\end{description}
				
					Longhand report by George Watt, in Church History Library, CR 100 317 (Historian's Office reports of speeches, 1845--1885), Brigham Young, 12 February 1854, transcribed by Bryce Gessell; view online \href{https://catalog.churchofjesuschrist.org/assets/1e4f947a-0644-47b0-b058-5dd134a82a00/0/0}{here}
					% brief description of original source material, with reference to JSP or somewhere else for more info
					% Background material: it might be helpful to have a note that says whether a given document was presented as a speech, was written in a journal, etc.
					% Also a comment here about how reliable the source might be, or what shape it is in, if there are large omissions or parts that are hard to understand, and so on.
					% Also a comment about when I transcribed it, whether I've checked it more than once, and so on
				
				\begin{description}
					\item[Text]
				\end{description}
				
					When I contemplate the subject of salvation, and rise \sout{up} before a congregation to speak upon that <all important matter,> \sout{subject} it has been but a few times in my life that I could see a bigining point to it, or a stoping place. There is such a multiplicity of principles, and sercumstances all interwoven so closely that it seems to be one eternity. I supose this is the case in reality. To be saved;---to be redeemed, and to have a right to the Celestial kingdom of God is to be intimately connected with the principles of eternity.
		
					I recolect once when I was preaching a question was asked me, ``What is \sout{the} Preisthood.'' The answer was ready, and perfectly simple in its nature, \sout{and} plain to be understood, <and> \sout{\&} \sout{is} couched in a a short sentence (viz) Preisthood is a perfect system of Government that rules and reigns in eternity.
					
					A question at once arises in the mind; ``\uline{where} \uline{is} \uline{eternity}.'' The answer is at hand, \uline{Eternity} is here;---we are in eternity just as much so as any other \sout{created} beings in heaven or on earth, Heavenly beings are no more in eternity than we are. We are in the midst of eternity; and when we become aquainted with the system of Gov\supers{t.} and the laws that rule in eternity, we shall then know it callculated to endure; to govern and control all things which are in heaven, and on the earth.
					
					The eternal preisthood of God,---the Gov\supers{t.} of God,---the laws of eternity, <is> \sout{are} a pure and perfect system of Gov\supers{t.}
		
		% -----
		\subsubsection{Instruction, 27 June 1854}
		% -----
			\label{EMS-1854-BY-27-JUN-1854}
			% \label{EMS-18XX-BY-DAY-3LETTERMONTH-YEAR}
			
			\begin{description}
					\item[Print Sources]
				\end{description}

					\begin{tabular}{p{0.5in}p{1in}p{0.5in}p{1in}}
					CDBY	& 		&  		&  \\
					\end{tabular}
				
				\begin{description}
					\item[Online Access] \href{https://catalog.churchofjesuschrist.org/assets/1dc08736-40cc-4a4d-86be-759370da1731/0/0}{Here}
					% \item[Analysis] \hyperref[XX]{Here}
				\end{description}
				
				\begin{description}
					\item[Background]
				\end{description}
				
					\phantom{.}
				
				\begin{description}
					\item[Original Source]
				\end{description}
				
					Church History Library, CR 100 318 (Historian's Office general Church minutes, 1839--1877), Brigham Young, 27 June 1854; view online \href{https://catalog.churchofjesuschrist.org/assets/1dc08736-40cc-4a4d-86be-759370da1731/0/0}{here}
					% brief description of original source material, with reference to JSP or somewhere else for more info
					% Background material: it might be helpful to have a note that says whether a given document was presented as a speech, was written in a journal, etc.
					% Also a comment here about how reliable the source might be, or what shape it is in, if there are large omissions or parts that are hard to understand, and so on.
					% Also a comment about when I transcribed it, whether I've checked it more than once, and so on
				
				\begin{description}
					\item[Text]
				\end{description}
				
					The Priesthood is a pure system of government which governs the world and a small portion has come down to earth. It is the constitutional laws and pure system by which men are to be governed.
					
		% -----
		\subsubsection{Instruction, 19 November 1854}
		% -----
			\label{EMS-1854-BY-19-NOV-1854}
			% \label{EMS-18XX-BY-DAY-3LETTERMONTH-YEAR}
			
			\begin{description}
					\item[Print Sources]
				\end{description}

					\begin{tabular}{p{0.5in}p{1in}p{0.5in}p{1in}}
					CDBY	& 		&  		&  \\
					\end{tabular}
				
				\begin{description}
					\item[Online Access] \href{https://catalog.churchofjesuschrist.org/assets/c05f1725-7304-4a19-ae73-91fc46573e0e/0/0}{Here}
					% \item[Analysis] \hyperref[XX]{Here}
				\end{description}
				
				\begin{description}
					\item[Background]
				\end{description}
				
					\phantom{.}
				
				\begin{description}
					\item[Original Source]
				\end{description}
				
					Church History Library, CR 100 317 (Historian's Office reports of speeches, 1845--1885), Brigham Young, 19 November 1854, transcribed by Bryce Gessell; view online \href{https://catalog.churchofjesuschrist.org/assets/c05f1725-7304-4a19-ae73-91fc46573e0e/0/0}{here}
					% brief description of original source material, with reference to JSP or somewhere else for more info
					% Background material: it might be helpful to have a note that says whether a given document was presented as a speech, was written in a journal, etc.
					% Also a comment here about how reliable the source might be, or what shape it is in, if there are large omissions or parts that are hard to understand, and so on.
					% Also a comment about when I transcribed it, whether I've checked it more than once, and so on
				
				\begin{description}
					\item[Text]
				\end{description}
				
					There is no kingdom without a space, and there is no space without a kingdom says the prophet Joseph. Now let me tell you there is no world made without that world has motion, and it is governed and controled by the law of the eternal preisthood; and if you want to know what preisthood is, it is a perfect system of Government by which planets and all creations are governed. Philosephers sit down, and tell us how the Gods came into existance, and how the first creation was made, and the planetts, it proves to me they expose their poor ignorance before all intelligence, and plants themselves in the pages of forgetfulness. There is no plannets nor kingdom, that mortal man can \sout{find} \sout{out} decifer and find out without the revelations of the Almighty. But I am not an asstronomer nor a philosopher, nor do I pofess any learning; I have none only what the Lord taught me, and what I got while I was in the woods rolling loggs, and while I was in the willderness; but there is knowledge and truth; and it is just as easy for the Lord to reveal to a person a mystery with regard to the plannets motion on their axes, that they can tell it to a congregation or to the world in one minute, as to be six monnths about it; just exactly, I want something that will do us \sout{some} good.
					
		% -----
		\subsubsection{Instruction, 3 December 1854}
		% -----
			\label{EMS-1854-BY-3-DEC-1854}
			% \label{EMS-18XX-BY-DAY-3LETTERMONTH-YEAR}
			
			\begin{description}
					\item[Print Sources]
				\end{description}

					\begin{tabular}{p{0.5in}p{1in}p{0.5in}p{1in}}
					CDBY	& 		&  		&  \\
					\end{tabular}
				
				\begin{description}
					\item[Online Access] \href{https://jod.mrm.org/2/136}{Here}
					% \item[Analysis] \hyperref[XX]{Here}
				\end{description}
				
				\begin{description}
					\item[Background]
				\end{description}
				
					\phantom{.}
				
				\begin{description}
					\item[Original Source]
				\end{description}
				
					I haven't been able to find any better source for this sermon than the JD account.
					% brief description of original source material, with reference to JSP or somewhere else for more info
					% Background material: it might be helpful to have a note that says whether a given document was presented as a speech, was written in a journal, etc.
					% Also a comment here about how reliable the source might be, or what shape it is in, if there are large omissions or parts that are hard to understand, and so on.
					% Also a comment about when I transcribed it, whether I've checked it more than once, and so on
				
				\begin{description}
					\item[Text]
				\end{description}
				
					The Spirit of the Lord, in teaching the people, in opening their minds to the principles of truth, does not infringe upon the laws God has given to mankind for their government; consequently, when the Lord made man, He made him an agent accountable to his God, with liberty to act and to do as he pleases, to a certain extent, in order to prove himself. There is a law that governs man thus far; but the law of the celestial kingdom, as I have frequently told you, is, and always will be, the same to all the children of Adam. When we talk of the celestial law which is revealed from heaven, that is, the Priesthood, we are talking about the principle of salvation, a perfect system of government, of laws and ordinances, by which we can be prepared to pass from one gate to another, and from one sentinel to another, until we go into the presence of our Father and God. This law has not always been upon the earth; and in its absence, other laws have been given to the children of men for their improvement, for their education, for their government, and to prove what they would do when left to control themselves; and what we now call tradition has grown out of these circumstances.
		
					There is so much of this, that I hardly dare to commence talking about it. It would require a lengthy discourse upon this particular point. Suffice it to say, the Lord has not established laws by which I am compelled to have my shoes made in a certain style. He has never given a law to determine whether I shall have a square-toed boot or a peaked-toed boot; whether I shall have a coat with the waist just under my arms, and the skirts down to my heels; or whether I shall have a coat like the one I have on. Intelligence, to a certain extent, was bestowed both upon Saint and sinner, to use independently, aside from whether they have the law of the Priesthood or not, or whether they have ever heard of it or not. ``I put into you intelligence,'' saith the Lord, ``that you may know how to govern and control yourselves, and make yourselves comfortable and happy on the earth; and give unto you certain privileges to act upon as independently in your sphere as I do in the government of heaven.''
					
					No matter whether we are Jew or Gentile, as the two classes of people are called; though Gentile signifies disobedient people; no matter whether we believe in the Koran as firmly as we now believe in the Bible; no matter whether we have been educated by the Jews, the Gentiles, or the Hottentots; whether we serve the true and the living God, or a lifeless image, if we are honest before the God we serve.